% Page 030

\seite{30}

\abschnitt{Nachtrag}
\pstart
Am 9.\ November 1918 war der schreckliche\ob
Krieg zu Ende gegangen durch Revolution und\ob
Zusammenbruch im eigenen Lande, der Waffen\ohb
stillstand zwang unsere Heere, sich bis über\ob
den Rhein zurückzuziehen. Durch den Rückzug hatte\ob
das Dorf viel zu leiden. Jedoch wurden die kopflos\ob
heimkehrenden Württemberger von der Bevölkerung\ob
herzlich empfangen und gut bewirtet. Bald darauf\ob
erfolgte der Durchzug der feindlichen Armeen. —

Kurze Zeit später kehrten auch unsere ausge\ohb
sandten Krieger in das Heimatdorf zurück. 4 be\ohb
fanden sich noch in Gefangenschaft. Der letzte (Willi Adr.\unsicher)\ob
kehrte nach 7 jähriger Abwesenheit am 20.\ 3. 1920\ob
zurück. Leider haben wir in unserem Orte 7\ob
Gefallene zu beklagen. Ihnen hat das Dorf durch\ob
ein Kriegergedächtniskreuz ein äußeres, bleibendes\ob
Erinnerungszeichen geschaffen. Die Einweihung fand am\ob
20.\ Juni\unsicher{} 1920 statt. Zu Ehren der Heimkehrenden\ob
wurde am 19.\ Juli 1925 ein Kriegerheim\ohb
fest veranstaltet. Zu diesem Zwecke stellte\ob
die Gemeinde 1000 Mk zur Verfügung. Eine Mu\ohb
sikkapelle trug zur Verschönerung der Feier bei.\ob
Nach Begrüßung der Heimkehrenden erfolgte ein\ob
Umzug durch die Straßen. Auch die Schulkinder\ob
begrüßten die Helden durch Gedichte und Lieder.\ob
Lehrer Stoffels von hier gedachte in einer Anrede\ob
der Gefallenen und dankte der Gemeinde im\ob
Namen der Krieger für das schöne Fest.
\pend
